% Applied Soft Computing 论文LaTeX模板
\documentclass[review]{elsarticle}

\usepackage{lineno,hyperref}
\modulolinenumbers[5]

\usepackage{graphicx}
\usepackage{amsmath}
\usepackage{algorithm}
\usepackage{algorithmic}
\usepackage{booktabs}
\usepackage{subcaption}

\journal{Applied Soft Computing}

\bibliographystyle{elsarticle-num}

\begin{document}

\begin{frontmatter}

\title{Probabilistic Logistics Demand Forecasting using Graph Variational Autoencoder and Normalizing Flow}

\author[inst1]{Your Name\corref{cor1}}
\ead{your.email@university.edu}

\author[inst1]{Co-author Name}
\ead{coauthor@university.edu}

\cortext[cor1]{Corresponding author}
\address[inst1]{Department, University, City, Country}

\begin{abstract}
Accurate demand forecasting with uncertainty quantification is crucial for logistics optimization.
This paper proposes a novel probabilistic forecasting framework that integrates Graph Variational
Autoencoder (GraphVAE) for spatial dependency modeling and Normalizing Flow for uncertainty
quantification. The GraphVAE learns dynamic network structures from logistics data, while the
Normalizing Flow models complex probability distributions for prediction intervals. Experiments
on real-world logistics data demonstrate that our method achieves XX\% improvement over baseline
methods in terms of CRPS, while providing reliable uncertainty estimates for decision-making.
\end{abstract}

\begin{keyword}
Demand forecasting \sep Graph neural networks \sep Variational autoencoder \sep
Normalizing flow \sep Uncertainty quantification \sep Logistics optimization
\end{keyword}

\end{frontmatter}

\linenumbers

\section{Introduction}
\label{sec:introduction}

Logistics demand forecasting plays a critical role in supply chain optimization~\cite{ref1}.
Traditional point forecasting methods fail to capture the inherent uncertainty in demand patterns...

% TODO: 补充引言内容
% - 物流需求预测的重要性
% - 现有方法的局限性
% - 不确定性量化的必要性
% - 本文贡献

\section{Related Work}
\label{sec:related}

\subsection{Time Series Forecasting}
% LSTM, GRU, Transformer等方法

\subsection{Graph Neural Networks}
% GCN, GAT, GraphSAGE等

\subsection{Probabilistic Forecasting}
% 概率预测方法综述

\section{Methodology}
\label{sec:method}

\subsection{Problem Formulation}
给定历史需求序列 $\mathbf{X} = \{\mathbf{x}_1, \mathbf{x}_2, ..., \mathbf{x}_T\}$，
其中 $\mathbf{x}_t \in \mathbb{R}^N$ 表示 $t$ 时刻 $N$ 个网点的需求，
目标是预测未来 $H$ 个时间步的需求分布 $p(\mathbf{y}_{T+1:T+H} | \mathbf{X})$。

\subsection{Graph Variational Autoencoder}
GraphVAE模型学习网点间的空间依赖关系...

\begin{equation}
\mathcal{L}_{VAE} = \mathbb{E}_{q_{\phi}(z|X)}[\log p_{\theta}(X|z)] - D_{KL}(q_{\phi}(z|X) || p(z))
\end{equation}

\subsection{Normalizing Flow}
Normalizing Flow通过可逆变换建模复杂概率分布...

\begin{equation}
\log p(x) = \log p(f(x)) + \log \left| \det \frac{\partial f(x)}{\partial x} \right|
\end{equation}

\subsection{End-to-End Framework}
% 整体框架描述

\section{Experiments}
\label{sec:experiments}

\subsection{Dataset Description}
我们使用某物流公司4年的真实需求数据，包含87,780条记录...

\subsection{Baseline Methods}
\begin{itemize}
    \item LSTM~\cite{lstm}
    \item GRU~\cite{gru}
    \item Transformer~\cite{transformer}
    \item ARIMA~\cite{arima}
    \item Prophet~\cite{prophet}
\end{itemize}

\subsection{Evaluation Metrics}
\begin{itemize}
    \item MAE: Mean Absolute Error
    \item RMSE: Root Mean Square Error
    \item CRPS: Continuous Ranked Probability Score
\end{itemize}

\subsection{Results and Analysis}

\begin{table}[ht]
\centering
\caption{Model Performance Comparison}
\label{tab:results}
\begin{tabular}{lccc}
\toprule
\textbf{Model} & \textbf{MAE} & \textbf{RMSE} & \textbf{CRPS} \\
\midrule
LSTM & 0.1033 & 0.2975 & 0.0723 \\
GRU & 0.1024 & 0.2912 & 0.0717 \\
Transformer & 0.0987 & 0.2845 & 0.0691 \\
ARIMA & 0.1523 & 0.3214 & 0.1218 \\
Prophet & 0.1412 & 0.3087 & 0.1130 \\
\textbf{Proposed} & \textbf{0.0891} & \textbf{0.2654} & \textbf{0.0535} \\
\bottomrule
\end{tabular}
\end{table}

% TODO: 插入图表
\begin{figure}[ht]
\centering
\includegraphics[width=0.9\textwidth]{../outputs/figures/figure1_model_comparison.pdf}
\caption{Model performance comparison across different metrics.}
\label{fig:comparison}
\end{figure}

\subsection{Ablation Study}
为了验证各组件的有效性，我们进行了消融实验...

\section{Conclusion}
\label{sec:conclusion}

本文提出了一种基于GraphVAE和Normalizing Flow的概率需求预测方法...

% TODO: 补充结论
% - 主要贡献总结
% - 实验结果总结
% - 未来工作方向

\section*{Acknowledgments}
This work was supported by...

\bibliography{references}

\end{document}
